% from ega1-4-fr.tex
\oldpage[I]{125}
à chacun des $X'_\lambda\times_S Y'_\mu$ est encore une immersion~; il en est
par suite de même de $\alpha\times_S\beta$ en vertu de
(\hyperref[I.4.2.4-fr]{4.2.4}).

En second lieu, prouvons qu'on peut aussi supposer que $X$ et $Y$ sont
\emph{affines}. En effet, soit $(U_i)$ (resp. $(V_j)$) un recouvrement de $X$
(resp. $Y$) par des ouverts affines, et soit $X'_i$ (resp. $Y'_j$) la
restriction de $X'$ (resp. $Y'$) à $X'\cap U_i$ (resp. $Y'\cap V_j$), qui est
un sous-préschéma fermé de $U_i$ (resp. $V_j$)~; $U_i\times_S V_j$ s'identifie
à la restriction de $X\times_S Y$ à
$p^{-1}(U_i)\cap q^{-1}(V_j)$ (\hyperref[I.3.2.7-fr]{3.2.7})~; et de même, si
$p'$, $q'$ sont les projections de $X'\times_S Y'$,
$X'_i\times_S Y'_j$ s'identifie à la restriction de $X'\times_S Y'$ à
${p'}^{-1}(X'_i)\cap{q'}^{-1}(Y'_j)$. Posons
$\gamma=\alpha\times_S\beta$~; on a par définition
$p\circ\gamma=\alpha\circ p'$, $q\circ\gamma=\beta\circ q'$~; comme
$X'_i=\alpha^{-1}(U_i)$, $Y'_j=\beta^{-1}(V_j)$, on a aussi
${p'}^{-1}(X'_i)=\gamma^{-1}(p^{-1}(U_i))$,
${q'}^{-1}(Y'_j)=\gamma^{-1}(q^{-1}(V_j))$, d'où
\[
  {p'}^{-1}(X'_i)\cap{q'}^{-1}(Y'_j)
  =\gamma^{-1}\bigl(p^{-1}(U_i)\cap q^{-1}(V_j)\bigr)
  =\gamma^{-1}(U_i\times_S V_j)
\]
on conclut comme dans la première partie du raisonnement.

Supposons donc $X$, $Y$, $S$ affines, et soient $B$, $C$, $A$ leurs anneaux
respectifs. Alors $B$ et $C$ sont des $A$-algèbres, $X'$ et $Y'$ des schémas
affines dont les anneaux sont des algèbres quotients $B'$, $C'$ de $B$ et
$C$ respectivement. En outre, on a
$\alpha=({}^a\rho,\widetilde{\rho})$, $\beta=({}^a\sigma,\widetilde{\sigma})$,
où $\rho$ et $\sigma$ sont respectivement les homomorphismes canoniques
$B\to B'$, $C\to C'$ (\hyperref[I.1.7.3-fr]{1.7.3}). Cela étant, on sait que
$X\times_S Y$ (resp. $X'\times_S Y'$) est un schéma affine d'anneau
$B\otimes_A C$ (resp. $B'\otimes_A C'$), et
$\alpha\times_S\beta=({}^a\tau,\widetilde{\tau})$, où $\tau$ est
l'homomorphisme $\rho\otimes\sigma$ de $B\otimes_A C$ dans
$B'\otimes_A C'$ (\hyperref[I.3.2.2-fr]{3.2.2} et
\hyperref[I.3.2.3-fr]{3.2.3})~; comme cet homomorphisme est surjectif,
$\alpha\times_S\beta$ est bien une immersion. En outre, si $\mathfrak{b}$
(resp. $\mathfrak{c}$) est le noyau de $\rho$ (resp. $\sigma$), le noyau de
$\tau$ est $u(\mathfrak{b})+v(\mathfrak{c})$, où $u$ (resp. $v$) est
l'homomorphisme $b\mapsto b\otimes1$ (resp. $c\mapsto1\otimes c$). Comme
$p=({}^a u,\widetilde{u})$ et $q=({}^a v,\widetilde{v})$, ce noyau correspond,
dans le spectre premier de $B\otimes_A C$, à l'ensemble fermé
$p^{-1}(X')\cap q^{-1}(Y')$ (\hyperref[I.1.2.2.1-fr]{1.2.2.1} et
\hyperref[I.1.1.2-fr]{1.1.2 (iii)}), ce qui achève la démonstration.

\begin{corollary}[4.3.2]
\label{I.4.3.2-fr}
Si $f:X\to Y$ est une immersion (resp. une immersion ouverte, une immersion
fermée) et un $S$-morphisme, $f_{(S')}$ est une immersion (resp. une immersion
ouverte, une immersion fermée) pour toute extension $S'\to S$ du préschéma de
base.
\end{corollary}

\subsection{Image réciproque d'un sous-préschéma.}
\label{subsection:I.4.4-fr}

\begin{proposition}[4.4.1]
\label{I.4.4.1-fr}
Soient $f:X\to Y$ un morphisme, $Y'$ un sous-préschéma (resp. un
sous-préschéma fermé, un préschéma induit sur un ouvert) de $Y$, $j:Y'\to Y$
le morphisme d'injection. Alors la projection $p:X\times_Y Y'\to X$ est une
immersion (resp. une immersion fermée, une immersion ouverte)~; le
sous-préschéma de $X$ associé à $p$ a $f^{-1}(Y')$ pour espace sous-jacent~;
en outre, si $j'$ est le morphisme d'injection de ce sous-préschéma dans $X$,
pour qu'un morphisme $h:Z\to X$ soit tel que $f\circ h:Z\to Y$ soit majoré
par $j$, il faut et il suffit que $h$ soit majoré par $j'$.
\end{proposition}

Comme $p=1_X\times_Y j$ (\hyperref[I.3.3.4-fr]{3.3.4}), la première assertion
résulte de (\hyperref[I.4.3.1-fr]{4.3.1})~; la seconde est un cas particulier
de (\hyperref[I.3.5.10-fr]{3.5.10}) (où il faut échanger les rôles de $X$ et
de $Y'$). Enfin, si on a $f\circ h=j\circ h'$, où $h'$ est un morphisme
$Z\to Y'$, il résulte de la définition du produit que l'on a $h=p\circ u$,
où $u$ est un morphisme $Z\to X\times_Y Y'$, d'où la dernière assertion.

Nous dirons que le sous-préschéma de $X$ ainsi défini est l'\emph{image
réciproque} du sous-préschéma $Y'$ de $Y$ par le morphisme $f$, terminologie
qui s'accorde avec celle introduite
